% Generated from source physical page 58; do not hand-edit.
\HLABeginPage
  \node[anchor=base west,inner sep=0bp,outer sep=0bp] at (266.3400,30.6000) {{\HLAFontArial\fontsize{7.9800bp}{9.1770bp}\selectfont\color{HLAc000000}IEEE Std 1516.2-2025}};
  \node[anchor=base west,inner sep=0bp,outer sep=0bp] at (158.5797,40.6197) {{\HLAFontArial\fontsize{7.9800bp}{9.1770bp}\selectfont\color{HLAc000000}IEEE Standard for Modeling and Simulation (M\&S) High Level Architecture (HLA)—}};
  \node[anchor=base west,inner sep=0bp,outer sep=0bp] at (228.3600,50.5795) {{\HLAFontArial\fontsize{7.9800bp}{9.1770bp}\selectfont\color{HLAc000000}Object Model Template (OMT) Specification}};
  \node[anchor=base west,inner sep=0bp,outer sep=0bp] at (301.5000,758.1600) {{\HLAFontTimes\fontsize{9.9600bp}{11.4540bp}\selectfont\color{HLAc000000}57}};
  \node[anchor=base west,inner sep=0bp,outer sep=0bp] at (229.8000,768.8400) {{\HLAFontArial\fontsize{7.9800bp}{9.1770bp}\selectfont\color{HLAc000000}Copyright © 2025 IEEE. All rights reserved.}};
  \node[anchor=base west,inner sep=0bp,outer sep=0bp] at (90.0000,78.6600) {{\HLAFontTimes\fontsize{9.9600bp}{11.4540bp}\selectfont\color{HLAc000000}A logical time represents a single point on the HLA time axis. A logical time interval represents the distance}};
  \node[anchor=base west,inner sep=0bp,outer sep=0bp] at (90.0000,90.6598) {{\HLAFontTimes\fontsize{9.9600bp}{11.4540bp}\selectfont\color{HLAc000000}between two points on the HLA time axis. During a federation execution, logical time may play two roles.}};
  \node[anchor=base west,inner sep=0bp,outer sep=0bp] at (90.0000,102.6596) {{\HLAFontTimes\fontsize{9.9600bp}{11.4540bp}\selectfont\color{HLAc000000}Federates may associate themselves with points on the HLA time axis, and they may associate some of their}};
  \node[anchor=base west,inner sep=0bp,outer sep=0bp] at (90.0000,114.6594) {{\HLAFontTimes\fontsize{9.9600bp}{11.4540bp}\selectfont\color{HLAc000000}activities, such as updating an instance attribute’s value or sending an interaction, with points on the HLA}};
  \node[anchor=base west,inner sep=0bp,outer sep=0bp] at (90.0000,126.6592) {{\HLAFontTimes\fontsize{9.9600bp}{11.4540bp}\selectfont\color{HLAc000000}time axis. These points on the HLA time axis are referred to as time positions when associated with a}};
  \node[anchor=base west,inner sep=0bp,outer sep=0bp] at (89.9990,138.6591) {{\HLAFontTimes\fontsize{9.9600bp}{11.4540bp}\selectfont\color{HLAc000000}federate and timestamps when associated with their activities. As a federation execution progresses,}};
  \node[anchor=base west,inner sep=0bp,outer sep=0bp] at (89.9990,150.6589) {{\HLAFontTimes\fontsize{9.9600bp}{11.4540bp}\selectfont\color{HLAc000000}federates may advance their time positions along the HLA time axis. Logical time intervals are used to}};
  \node[anchor=base west,inner sep=0bp,outer sep=0bp] at (89.9990,162.6587) {{\HLAFontTimes\fontsize{9.9600bp}{11.4540bp}\selectfont\color{HLAc000000}specify a federate’s lookahead and by the RTI to perform arithmetic operations with logical time values.}};
  \node[anchor=base west,inner sep=0bp,outer sep=0bp] at (89.9980,184.6783) {{\HLAFontTimes\fontsize{9.9600bp}{11.4540bp}\selectfont\color{HLAc000000}The specific strategy used to advance time in a simulation is driven by the simulation’s purpose. For}};
  \node[anchor=base west,inner sep=0bp,outer sep=0bp] at (89.9980,196.6781) {{\HLAFontTimes\fontsize{9.9600bp}{11.4540bp}\selectfont\color{HLAc000000}example, faster-than-real time, event-stepped simulations are frequently used in analysis applications,}};
  \node[anchor=base west,inner sep=0bp,outer sep=0bp] at (89.9980,208.6779) {{\HLAFontTimes\fontsize{9.9600bp}{11.4540bp}\selectfont\color{HLAc000000}whereas real time, time-stepped simulations are often used in training applications. The strategy chosen for}};
  \node[anchor=base west,inner sep=0bp,outer sep=0bp] at (89.9980,220.6777) {{\HLAFontTimes\fontsize{9.9600bp}{11.4540bp}\selectfont\color{HLAc000000}advancing time is an important consideration in the design of a simulation, because it affects both simulation}};
  \node[anchor=base west,inner sep=0bp,outer sep=0bp] at (89.9970,232.6776) {{\HLAFontTimes\fontsize{9.9600bp}{11.4540bp}\selectfont\color{HLAc000000}performance and the ability of a simulation to interoperate with other simulations in a federation. For this}};
  \node[anchor=base west,inner sep=0bp,outer sep=0bp] at (89.9970,244.6774) {{\HLAFontTimes\fontsize{9.9600bp}{11.4540bp}\selectfont\color{HLAc000000}reason, it is important to document how a simulation supports the advancement of time in a federate’s SOM}};
  \node[anchor=base west,inner sep=0bp,outer sep=0bp] at (89.9970,256.6772) {{\HLAFontTimes\fontsize{9.9600bp}{11.4540bp}\selectfont\color{HLAc000000}and relevant SOM modules.}};
  \node[anchor=base west,inner sep=0bp,outer sep=0bp] at (89.9970,278.6370) {{\HLAFontTimes\fontsize{9.9600bp}{11.4540bp}\selectfont\color{HLAc000000}IEEE Std 1516.1-2025 provides a variety of time management services to control the advancement of}};
  \node[anchor=base west,inner sep=0bp,outer sep=0bp] at (89.9970,290.6368) {{\HLAFontTimes\fontsize{9.9600bp}{11.4540bp}\selectfont\color{HLAc000000}federates along the HLA time axis. These services allow multiple federates with differing internal time}};
  \node[anchor=base west,inner sep=0bp,outer sep=0bp] at (89.9970,302.6366) {{\HLAFontTimes\fontsize{9.9600bp}{11.4540bp}\selectfont\color{HLAc000000}advancement strategies to interoperate in a meaningful way. For this reason, it is vital for all federation}};
  \node[anchor=base west,inner sep=0bp,outer sep=0bp] at (89.9970,314.6364) {{\HLAFontTimes\fontsize{9.9600bp}{11.4540bp}\selectfont\color{HLAc000000}members to agree on how logical times and logical time intervals are represented across the federation, and}};
  \node[anchor=base west,inner sep=0bp,outer sep=0bp] at (89.9970,326.6362) {{\HLAFontTimes\fontsize{9.9600bp}{11.4540bp}\selectfont\color{HLAc000000}document this agreement in the federation's FOM and relevant FOM modules.}};
  \node[anchor=base west,inner sep=0bp,outer sep=0bp] at (89.9970,348.6558) {{\HLAFontTimes\fontsize{9.9600bp}{11.4540bp}\selectfont\color{HLAc000000}During federation execution, logical times and logical time intervals are represented as instances of the time}};
  \node[anchor=base west,inner sep=0bp,outer sep=0bp] at (89.9980,360.6556) {{\HLAFontTimes\fontsize{9.9600bp}{11.4540bp}\selectfont\color{HLAc000000}representation abstract data type (TRADT). This abstract data type is provided to the RTI when federation}};
  \node[anchor=base west,inner sep=0bp,outer sep=0bp] at (89.9980,372.6554) {{\HLAFontTimes\fontsize{9.9600bp}{11.4540bp}\selectfont\color{HLAc000000}execution is created so that an RTI knows how to use the logical times and logical time intervals. This table}};
  \node[anchor=base west,inner sep=0bp,outer sep=0bp] at (89.9980,384.6552) {{\HLAFontTimes\fontsize{9.9600bp}{11.4540bp}\selectfont\color{HLAc000000}allows the data type and semantics of the abstract data type for logical times and logical time intervals to be}};
  \node[anchor=base west,inner sep=0bp,outer sep=0bp] at (89.9980,396.6551) {{\HLAFontTimes\fontsize{9.9600bp}{11.4540bp}\selectfont\color{HLAc000000}explained in FOMs, FOM modules, SOMs, and SOM modules.}};
  \node[anchor=base west,inner sep=0bp,outer sep=0bp] at (89.9980,418.6746) {{\HLAFontTimes\fontsize{9.9600bp}{11.4540bp}\selectfont\color{HLAc000000}It is also important to define the lookahead characteristics of both federates and federations. At the}};
  \node[anchor=base west,inner sep=0bp,outer sep=0bp] at (89.9980,430.6744) {{\HLAFontTimes\fontsize{9.9600bp}{11.4540bp}\selectfont\color{HLAc000000}individual federate level, the specific means by which lookahead is supported may be a factor in assessing}};
  \node[anchor=base west,inner sep=0bp,outer sep=0bp] at (89.9980,442.6743) {{\HLAFontTimes\fontsize{9.9600bp}{11.4540bp}\selectfont\color{HLAc000000}compatibility with other potential federates for a given federation application. For federations, it is important}};
  \node[anchor=base west,inner sep=0bp,outer sep=0bp] at (89.9980,454.6741) {{\HLAFontTimes\fontsize{9.9600bp}{11.4540bp}\selectfont\color{HLAc000000}to document a common set of characteristics for how lookahead will be supported across the full federation.}};
  \node[anchor=base west,inner sep=0bp,outer sep=0bp] at (89.9980,466.6739) {{\HLAFontTimes\fontsize{9.9600bp}{11.4540bp}\selectfont\color{HLAc000000}Thus, SOMs, SOM modules, FOMs and FOM modules all include representations of lookahead as well as}};
  \node[anchor=base west,inner sep=0bp,outer sep=0bp] at (89.9980,478.6737) {{\HLAFontTimes\fontsize{9.9600bp}{11.4540bp}\selectfont\color{HLAc000000}timestamps.}};
  \node[anchor=base west,inner sep=0bp,outer sep=0bp] at (89.9980,500.6933) {{\HLAFontTimes\fontsize{9.9600bp}{11.4540bp}\selectfont\color{HLAc000000}An abstract data type for lookahead is supplied when a federation execution is created. The lookahead entry}};
  \node[anchor=base west,inner sep=0bp,outer sep=0bp] at (89.9980,512.6931) {{\HLAFontTimes\fontsize{9.9600bp}{11.4540bp}\selectfont\color{HLAc000000}in this table allows the data type and semantics of this abstract data type to be documented in FOMs, SOMs,}};
  \node[anchor=base west,inner sep=0bp,outer sep=0bp] at (89.9980,524.6929) {{\HLAFontTimes\fontsize{9.9600bp}{11.4540bp}\selectfont\color{HLAc000000}FOM modules, and SOM modules.}};
  \node[anchor=base west,inner sep=0bp,outer sep=0bp] at (89.9980,546.6527) {{\HLAFontTimes\fontsize{9.9600bp}{11.4540bp}\selectfont\color{HLAc000000}Semantic differences exist between the way time is represented for the purpose of depicting timestamps}};
  \node[anchor=base west,inner sep=0bp,outer sep=0bp] at (89.9980,558.6525) {{\HLAFontTimes\fontsize{9.9600bp}{11.4540bp}\selectfont\color{HLAc000000}versus calculating lookahead. When depicting timestamps, time can be considered to be an absolute value on}};
  \node[anchor=base west,inner sep=0bp,outer sep=0bp] at (89.9980,570.6523) {{\HLAFontTimes\fontsize{9.9600bp}{11.4540bp}\selectfont\color{HLAc000000}the HLA time axis, and thus, time comparisons can be drawn to determine if one timestamp is greater than}};
  \node[anchor=base west,inner sep=0bp,outer sep=0bp] at (89.9990,582.6521) {{\HLAFontTimes\fontsize{9.9600bp}{11.4540bp}\selectfont\color{HLAc000000}another. Lookahead, in contrast, represents a duration of time, which can be added to timestamps but is}};
  \node[anchor=base west,inner sep=0bp,outer sep=0bp] at (89.9990,594.6519) {{\HLAFontTimes\fontsize{9.9600bp}{11.4540bp}\selectfont\color{HLAc000000}generally not used for comparison purposes.}};
  \node[anchor=base west,inner sep=0bp,outer sep=0bp] at (89.9990,616.6715) {{\HLAFontTimes\fontsize{9.9600bp}{11.4540bp}\selectfont\color{HLAc000000}Two time representations, HLAinteger64Time and HLAfloat64Time, are required to be supported by the}};
  \node[anchor=base west,inner sep=0bp,outer sep=0bp] at (89.9990,628.6713) {{\HLAFontTimes\fontsize{9.9600bp}{11.4540bp}\selectfont\color{HLAc000000}HLA and are described in detail in IEEE Std 1516.1-2025. These types shall be provided by all RTIs.}};
  \node[anchor=base west,inner sep=0bp,outer sep=0bp] at (89.9990,640.6711) {{\HLAFontTimes\fontsize{9.9600bp}{11.4540bp}\selectfont\color{HLAc000000}Federations may choose to use these default data types, supply their own time representations, or not supply}};
  \node[anchor=base west,inner sep=0bp,outer sep=0bp] at (89.9990,652.6709) {{\HLAFontTimes\fontsize{9.9600bp}{11.4540bp}\selectfont\color{HLAc000000}any time representation (e.g., NA).}};
  \node[anchor=base west,inner sep=0bp,outer sep=0bp] at (89.9990,674.6904) {{\HLAFontArial\fontsize{9.9600bp}{11.4540bp}\selectfont\bfseries\color{HLAc000000}4.8.2 Table format}};
  \node[anchor=base west,inner sep=0bp,outer sep=0bp] at (89.9990,696.6503) {{\HLAFontTimes\fontsize{9.9600bp}{11.4540bp}\selectfont\color{HLAc000000}The template that shall be used for recording time representation is illustrated in Table 16.}};
\HLAEndPage
